\documentclass[10pt]{article}
\usepackage[a4paper,margin=25mm]{geometry}
\usepackage[T1]{fontenc}
\usepackage{lmodern,amsmath,amssymb,amsthm,booktabs,array}
\usepackage[hyphens]{url}
\usepackage[colorlinks=true,linkcolor=blue,citecolor=blue,urlcolor=blue]{hyperref}
\newtheorem{theorem}{Theorem}
\newtheorem{lemma}[theorem]{Lemma}
\newtheorem{corollary}[theorem]{Corollary}
\newtheorem{proposition}[theorem]{Proposition}
\newcommand{\gam}{\gamma}
\title{Weighted toroidal extensions of Liang (2025):\\
OEIS A394697, history revision 21,\\6 May 2026 conjecture}
\author{Haobo Ma\thanks{ChronoAI Pte Ltd, Singapore.}\and
Wenlin Zhang\thanks{National University of Singapore; ChronoAI Pte Ltd, Singapore.}}
\date{7 November 2026}
\begin{document}
\maketitle
\begin{abstract}
Extending Liang's (2025) free-board peak-entropy argument, we prove
Kot\v{e}\v{s}ovec's conjecture in OEIS A394697, history revision 21
(6 May 2026; entry snapshot revision 24): the largest coefficient of the independence
polynomial of the $n\times n$ toroidal king graph has $n^2$th root tending
to the square-tiling constant A247413. The limit holds through all integers
$n\geq3$. An empty final row and column compare toroidal counts with free
counts; a maximum-versus-sum estimate then transfers their entropy to the
largest coefficient. The same
argument applies at every fixed positive activity and to rectangular
boards with free, cylindrical or toroidal boundaries, with explicit
boundary-to-area error bounds. Positive moments of the number of kings
have the same leading entropy. We also disprove the entry's scalar
log-concavity comment and show that such failures occur infinitely often.
\end{abstract}

\section{The question and its normalization}
At history revision 21 of A394697~\cite{a394697} (6 May 2026),
Kot\v{e}\v{s}ovec conjectured that $a(n)^{1/n^2}$ ``tends to A247413.''
The conjecture remains in entry revision 24. It asserts convergence of a modal
coefficient, rather than of the total number of independent sets. We prove
that assertion with the literal normalization and constant given in the
entry. In particular, no restriction to even board sizes is necessary.

For $n\geq3$, let $G_n$ have vertex set $(\mathbb Z/n\mathbb Z)^2$.
Distinct vertices are adjacent if their cyclic distances in both
coordinates are at most one. Write
\[
 I_n(x)=\sum_{k=0}^{n^2}c_{n,k}x^k,\qquad
 T_n=I_n(1),\qquad a_n=\max_{0\leq k\leq n^2}c_{n,k}.
\]
Thus $c_{n,k}$ counts independent sets of size $k$, including the empty
set when $k=0$; zero coefficients remain in the range of the maximum.
Let $F_n$ count independent sets on the free $n\times n$ king board,
where ordinary coordinate differences replace cyclic distances, and set
$F_0=1$. Let $B_m$ count tilings of an $m\times m$ square by unit and
$2\times2$ squares, including $B_0=1$. The constant in the conjecture is
\[
 \gam=\lim_{m\to\infty}B_m^{1/m^2},
 \tag{1}\label{constant}
\]
the constant recorded in A247413~\cite{a247413}; the tiling sequence is
A063443~\cite{a063443}. Its existence is included in the proof below.

\begin{theorem}\label{main}
The limit in \eqref{constant} exists, and
\[
 \lim_{n\to\infty}a_n^{1/n^2}
 =\lim_{n\to\infty}T_n^{1/n^2}
 =\lim_{n\to\infty}F_n^{1/n^2}=\gam.
\]
Equivalently, for every $\varepsilon>0$ there is an integer $N\geq3$
such that $|a_n^{1/n^2}-\gam|<\varepsilon$ for every integer $n\geq N$.
\end{theorem}

Free-board entropy predates this conjecture. Calkin, James, Purvis, Race,
Schneider and Yancey~\cite{kings}, Sections 3--4, study free king counts
and bound their exponential growth. Liang~\cite{liang}, Proposition 3.1,
identifies the common leading entropy of free-rectangle total counts and
largest coefficients using the number of possible cardinalities. His
later study~\cite{liang2}, Sections 3.2 and 4, compares king-board growth
with boundary-dependent corrections in related grid models. The present
proof combines that maximum-versus-sum mechanism with a direct
two-coordinate seam comparison, keeping the tiling side shift explicit.
Hildebrant~\cite{hildebrant}, Proposition 2.5 and Remark 2.6, supplies the transfer-operator
trace formulation used in the A394697 entry; it is useful for finite
examples but is not needed for the limiting argument.

\section{Free entropy, empty seams, and the peak}
\begin{lemma}\label{entropy}
The sequence $n^{-2}\log F_n$ converges to
\[
 h=\inf_{k\geq1}\frac{\log F_k}{k^2},\qquad
 \frac{\log2}{4}\leq h\leq\log2.
\]
\end{lemma}
\begin{proof}
Fix $k\geq1$ and write $n=qk+r$, where $0\leq r<k$.
Partition the square into $q^2$ disjoint $k\times k$ blocks and
$n^2-q^2k^2$ remaining vertices. Restricting an independent set to
these parts is injective. Each block has at most $F_k$ choices;
each remaining vertex has at most two. Edges between parts only
remove choices. Hence
\[
 F_n\leq F_k^{q^2}\,2^{n^2-q^2k^2}.
\]
Divide its logarithm by $n^2$ and let $n\to\infty$ with $k$ fixed.
Since $q/n\to1/k$ and $(n^2-q^2k^2)/n^2\to0$, the limsup is at
most $\log F_k/k^2$. Taking the infimum over $k$ gives a limsup
at most $h$; the definition of $h$ bounds every term below by $h$.
This proves convergence. All subsets of the vertices whose two
coordinates are even are independent, so
$F_n\geq2^{\lceil n/2\rceil^2}$. Together with $F_n\leq2^{n^2}$,
this gives the stated bounds on $h$.
\end{proof}

\begin{proof}[Proof of Theorem~\ref{main}]
First compare boundary conditions. Embed a free $(n-1)\times(n-1)$
configuration in coordinates $0,\ldots,n-2$ of the $n$-torus and
leave its final row and column empty. For two occupied coordinates,
the ordinary gap is at most $n-2$. Therefore a cyclic gap at most
one cannot arise from a wraparound gap of $n-1$: it is exactly an
ordinary gap at most one. This is a size-preserving injection into
toroidal independent sets. Conversely, forgetting the wraparound
edges preserves every toroidal independent set as a free one. Thus
\[
 F_{n-1}\leq T_n\leq F_n\qquad(n\geq3).
 \tag{2}\label{seam}
\]
Both bounding logarithms divided by $n^2$ tend to $h$, because
\[
 \frac{\log F_{n-1}}{n^2}
 =\left(\frac{n-1}{n}\right)^2
   \frac{\log F_{n-1}}{(n-1)^2}\longrightarrow h.
\]
Consequently $n^{-2}\log T_n\to h$.

There are $n^2+1$ cardinality classes, including zero. Nonnegativity
and summation of the coefficients give
\[
 a_n\leq T_n\leq(n^2+1)a_n.
 \tag{3}\label{peak}
\]
Since the empty set ensures $a_n\geq1$, logarithms are defined, and
\[
 \frac{\log T_n-\log(n^2+1)}{n^2}
 \leq\frac{\log a_n}{n^2}\leq\frac{\log T_n}{n^2}.
\]
The loss $\log(n^2+1)/n^2$ tends to zero. Thus
$n^{-2}\log a_n\to h$.

Finally, encode a $2\times2$ tile in an $m\times m$ square by its
lower-left cell, in $\{0,\ldots,m-2\}^2$. Two such tiles overlap
exactly when the lower-left cells differ by at most one in both
coordinates. Disjoint tiles therefore correspond exactly to a free
king configuration on side $m-1$. Every other cell has a unique unit
tile, giving the bijection
\[
 B_m=F_{m-1}\qquad(m\geq1).
 \tag{4}\label{tiling}
\]
It follows that $m^{-2}\log B_m\to h$ as well. Exponentiating all
these limits identifies \eqref{constant} with $e^h$ and proves every
asserted limit. This also establishes the displayed $\varepsilon$--$N$
statement through all integers, without a parity condition.
\end{proof}

\section{Weighted counts and a scalar counterexample}
For $z>0$, let $F_n(z)$ and $C_n(z)$ be the independent-set
polynomials evaluated at $z$ for the free board and the board with
horizontal periodicity only. Set $T_n(z)=I_n(z)$ and
$A_n(z)=\max_k c_{n,k}z^k$.

\begin{corollary}\label{weighted}
For each fixed real $z>0$ there is a number $h(z)$ such that
\[
 \lim_{n\to\infty}\frac{\log F_n(z)}{n^2}
 =\lim_{n\to\infty}\frac{\log C_n(z)}{n^2}
 =\lim_{n\to\infty}\frac{\log T_n(z)}{n^2}
 =\lim_{n\to\infty}\frac{\log A_n(z)}{n^2}=h(z).
\]
Moreover $\frac14\log(1+z)\leq h(z)\leq\log(1+z)$.
In particular the total-count sequences A067958 and A212269 have
$n^2$th-root limit $\gam$.
\end{corollary}
\begin{proof}
All configurations have positive weight $z^{|S|}$. Restriction to
disjoint pieces factors this weight, so Lemma~\ref{entropy} holds
with $F_n(z)$ and with $1+z$ replacing $2$. Subsets of the
even-coordinate vertices supply the lower bound. The injections in
\eqref{seam} preserve $|S|$, and the edge sets are nested, giving
\[
 F_{n-1}(z)\leq T_n(z)\leq C_n(z)\leq F_n(z).
\]
The same squeeze applies. Finally
$A_n(z)\leq T_n(z)\leq(n^2+1)A_n(z)$, so \eqref{peak} applies
to weighted coefficients. At $z=1$, A067958 counts the toroidal
configurations and A212269 the cylindrical ones~\cite{boundaries}.
\end{proof}

\begin{corollary}[Rectangles and explicit entropy bounds]\label{rectangles}
Let $m,n\geq3$, $z>0$, and let $Z^b_{m,n}(z)$ be the partition
function on the $m\times n$ king board with boundary condition
$b\in\{\mathrm{free},\mathrm{cylinder},\mathrm{torus}\}$. In the
cylindrical case the width $m$ is periodic. Let $P^b_{m,n}(z)$ be its
largest weighted coefficient and put $L=\log(1+z)$. Then
\[
 \frac{\log F_{m,n}(z)}{(m+1)(n+1)}\leq h(z)
 \leq\frac{\log F_{m,n}(z)}{mn},
 \tag{5}\label{bracket}
\]
where $F_{m,n}=Z^{\mathrm{free}}_{m,n}$, and
\[
 -\frac{(m+n-1)L+\log(mn+1)}{mn}
 \leq \frac{\log P^b_{m,n}(z)}{mn}-h(z)
 \leq\frac{(m+n+1)L}{mn}.
 \tag{6}\label{error}
\]
Consequently all three partition functions and their largest weighted
coefficients have entropy $h(z)$ whenever $\min(m,n)\to\infty$,
without a restriction on aspect ratio.
\end{corollary}
\begin{proof}
Partition a large free square into $m\times n$ rectangles and a
remainder. Restriction multiplies the rectangular weights, and the
remainder contributes at most $(1+z)$ per vertex. Taking its entropy
limit gives the upper bound in \eqref{bracket}. Conversely, place
independently chosen $m\times n$ configurations in a large square,
with an empty row and column between successive rectangles. Their
density is $1/((m+1)(n+1))$; this gives the lower bound. This is the
separated-block method used for king entropy by Calkin et al.~\cite{kings},
Lemma 4; Liang~\cite{liang2}, Section 2.2, uses analogous brackets
for the ordinary grid model.

Write $s=m+n-1$. Restricting a free configuration to its
$(m-1)\times(n-1)$ subrectangle leaves $s$ vertices, so
\[
 F_{m,n}(z)\leq(1+z)^sF_{m-1,n-1}(z),\qquad
 F_{m-1,n-1}(z)\leq Z^b_{m,n}(z)\leq F_{m,n}(z).
\]
The second inequality uses empty final coordinate lines, just as
in \eqref{seam}. Thus $\log Z^b_{m,n}/(mn)\geq h(z)-sL/(mn)$.
The lower bracket also gives
$\log F_{m,n}/(mn)\leq h(z)+(m+n+1)h(z)/(mn)$;
use $h(z)\leq L$. Finally
$P^b_{m,n}\leq Z^b_{m,n}\leq(mn+1)P^b_{m,n}$ gives \eqref{error}.
Every displayed error tends to zero when both side lengths grow.
\end{proof}

\begin{corollary}[Marked configurations]\label{marked}
For fixed $z>0$ and $p\geq0$, write
$Z^b_{m,n}(z)=\sum_k c^b_{m,n,k}z^k$ and define
\[
 M^b_{m,n,p}(z)=\sum_{k\geq1}k^p c^b_{m,n,k}z^k.
\]
Then $\log M^b_{m,n,p}(z)/(mn)\to h(z)$ as
$\min(m,n)\to\infty$, for each of the three boundaries.
For $p=1$, this counts configurations with one distinguished king.
\end{corollary}
\begin{proof}
Since $1\leq k^p\leq(mn)^p$ for $1\leq k\leq mn$,
$Z^b_{m,n}(z)-1\leq M^b_{m,n,p}(z)\leq(mn)^p Z^b_{m,n}(z)$.
Corollary~\ref{rectangles} and $h(z)\geq\frac14\log(1+z)>0$
show that subtracting $1$ does not change the entropy. The factor
$(mn)^p$ contributes $p\log(mn)/(mn)\to0$.
The free-board $p=1,z=1$ case reproves Liang~\cite{liang},
Proposition 4.1; the statement here includes periodic boundaries.
\end{proof}

\begin{corollary}[A189889: even sizes and leading density]\label{packing}
The degree $\alpha_n$ of $I_n$ satisfies
$\lfloor n/2\rfloor^2\leq\alpha_n\leq\lfloor n^2/4\rfloor$.
In particular $\alpha_n=n^2/4$ for even $n\geq4$, and
$\alpha_n/n^2\to1/4$.
\end{corollary}
\begin{proof}
Choose both coordinates from $0,2,\ldots,2\lfloor n/2\rfloor-2$
to obtain the lower bound. Each cyclic $2\times2$ square contains
at most one king. Summing this inequality over all $n^2$ such
squares counts each king four times, giving $4\alpha_n\leq n^2$.
The bounds coincide at even $n$ and have the same limiting density.
A189889~\cite{capacity} records the stronger exact formula
$\lfloor n\lfloor n/2\rfloor/2\rfloor$, attributed there to
Watkins and Ricci (2004); only its even branch and leading density
are reproved here.
\end{proof}

\begin{proposition}\label{logconcavity}
The scalar sequence $(a_n)_{n\geq3}$ is not log-concave:
$a_4^2=3136<5400=a_3a_5$. In fact the inequality
$a_n^2<a_{n-1}a_{n+1}$ holds for infinitely many $n$.
\end{proposition}
\begin{proof}
The entry gives $a_3=9$, $a_4=56$, $a_5=600$; the same values
follow from the finite polynomials below. For the infinite assertion,
suppose instead that $a_n^2\geq a_{n-1}a_{n+1}$ for every sufficiently
large $n$. Then the successive differences of $b_n=\log a_n$ are
eventually nonincreasing, hence bounded above by a constant $D$.
Telescoping gives $b_n\leq E+Dn$ for a fixed $E$. Since $b_n\geq0$,
this forces $b_n/n^2\to0$, contradicting
Theorem~\ref{main} and $h\geq\log2/4>0$.
\end{proof}

This addresses the literal comment ``a(n) is empirically log-concave
for n $\leq$ 21'' in A394697. It is
distinct from log-concavity of the coefficient list $(c_{n,k})_k$
of a single independence polynomial.

\section{Coverage and finite examples}
The source locators below refer to the OEIS data at revision
\begin{center}\small\texttt{77073cc8105b8042be07ec0d885da5ea5960ff3c}.\end{center}
\begin{center}\small
\begin{tabular}{@{}>{\raggedright\arraybackslash}p{.42\linewidth}>{\raggedright\arraybackslash}p{.52\linewidth}@{}}
\toprule
Source item & Mathematical coverage\\
\midrule
A394697, history rev. 21, 6 May 2026 conjecture & Proved: Theorem~\ref{main}, every integer $n\geq3$; extending Liang (2025).\\
A394697, scalar log-concavity comment & Refuted (explicit counterexample): Proposition~\ref{logconcavity}; infinitely many failures.\\
A394697, empirical increase of $n^{-2}\log a_n$ for $3\leq n\leq21$ & Verified on the listed values. Universal monotonicity is not asserted in the source; not covered here.\\
A394697, transfer trace and examples & Reproved; prior formula Hildebrant (2026), Prop. 2.5. Complete coefficients checked for $3\leq n\leq8$.\\
A063443, Woods's 27 Aug. 2011 tiling/free-board comment & Reproved; prior identity Woods (2011), bijection \eqref{tiling}.\\
A247413 and A063443, tiling entropy limit & Reproved; prior limit credited to McKay (1996) in A063443; Theorem~\ref{main}.\\
A067958, toroidal total count & Limit proved in Corollary~\ref{weighted}.\\
A212269, cylindrical limit formula & Reproved; prior formula Kot\v{e}\v{s}ovec, entry rev. 34 (2024)~\cite{boundaries}.\\
A189889, maximum number of kings & Reproved in part; even branch and density in Corollary~\ref{packing}. Exact odd branch not covered.\\
Liang (2025), Props. 3.1 and 4.1 & Reproved and extended: free peak and marked-count entropy; Corollaries~\ref{rectangles}--\ref{marked}.\\
Liang (2025), arXiv:2507.04007, Secs. 3.2, 4 & Prior boundary comparisons cited; convergence and values of surface coefficients not covered.\\
Weighted rectangular bounds and periodic marked moments & Not posed in these sources; proved in Corollaries~\ref{rectangles}--\ref{marked}.\\
\bottomrule
\end{tabular}
\end{center}

For completeness, the transfer trace in A394697 follows by indexing
rows by the independent subsets $U$ of the cycle $C_n$. Put
$M_n(x)_{U,V}=x^{|U|}$ when no vertex of $U$ equals or neighbors
a vertex of $V$, and put it equal to zero otherwise. Expanding
$\operatorname{tr}(M_n(x)^n)$ sums over ordered cyclic lists of
$n$ rows. Consecutive-row compatibility is exactly king independence,
including the last and first rows, and the product of entries is
$x^{\sum |U|}$. Each labelled toroidal configuration determines one
such list; no quotient by cyclic rotations is taken. Hence the trace
equals $I_n(x)$.

The complete polynomials at the first three sizes are
\begin{align*}
 I_3(x)&=1+9x,\\
 I_4(x)&=1+16x+56x^2+48x^3+12x^4,\\
 I_5(x)&=1+25x+200x^2+600x^3+575x^4+10x^5.
\end{align*}
The accompanying \texttt{verify.py} uses row-state transfer and,
independently for $n=3,4$, a vertex include/exclude recurrence for all
three boundary conditions. It reproduces complete toroidal
coefficient lists for $3\leq n\leq8$ and tests the weighted
inequalities at the exact rational activities $1/2,1,2$.
\begin{center}
\begin{tabular}{@{}rrrr@{}}
\toprule $n$ & $a_n$ & $T_n$ & $F_n$\\\midrule
3&9&10&35\\
4&56&133&314\\
5&600&1411&6427\\
6&14940&42938&202841\\
7&576856&1796859&12727570\\
8&42752576&157763829&1355115601\\\bottomrule
\end{tabular}
\end{center}
These finite examples are separate from the proofs of the infinite
limits.
Exact integer comparisons of the A394697 b-file values give
$a_n^2<a_{n-1}a_{n+1}$ for every $4\leq n\leq20$, and
$a_{n+1}^{n^2}>a_n^{(n+1)^2}$ for every $3\leq n\leq20$.
The latter verifies the entry's finite entropy-increase observation;
it does not imply monotonicity at all sizes. The accompanying
\texttt{sibling\_checks.py} also checks rectangular seam and
coefficient inequalities at $z=1/2,1,2$, for $3\leq m\leq6$ and
$3\leq n\leq8$, and compares the two enumeration methods on
$3\times3$, $3\times4$, and $4\times3$ boards.

\paragraph{Formal statements and written proofs.}
At source revision
\begin{center}\small\texttt{074df09ddbbef1843ddb383577bfaa827f753e36},\end{center}
the Lean declaration
\href{https://github.com/the-omega-institute/trureturing/blob/074df09ddbbef1843ddb383577bfaa827f753e36/D5/S3/Combinatorics/Graph/Pressure/ToroidalKingsPeakEntropy.lean}{\texttt{ToroidalKingsPeakEntropy.result}}
states the $\varepsilon$--$N$ limit in Theorem~\ref{main} for the
maximum over all cardinalities, with the constant defined from the
tiling sequence. In the same namespace hierarchy,
\href{https://github.com/the-omega-institute/trureturing/blob/074df09ddbbef1843ddb383577bfaa827f753e36/D5/S3/Combinatorics/Graph/Pressure/ToroidalKingsEmptySeam.lean}{\texttt{ToroidalKingsEmptySeam}}
contains \texttt{free\_le\_torus}, \texttt{torus\_le\_free},
\texttt{count\_le\_mul\_peak}, \texttt{tiling\_free\_identity},
and \texttt{free\_log\_limit}, corresponding to the seam, cardinality,
tiling, and free-entropy statements above. Corollary~\ref{weighted},
Corollaries~\ref{rectangles}--\ref{packing},
Proposition~\ref{logconcavity}, and the displayed trace argument have
the written proofs given here. The finite coefficient calculations
are specified in \texttt{checks.json} and are reproducible with
Python 3 using \texttt{verify.py}.
The supporting declaration
\href{https://github.com/the-omega-institute/trureturing/blob/074df09ddbbef1843ddb383577bfaa827f753e36/D5/S3/Combinatorics/Graph/Pressure/RectangularCountEntropy.lean}{\texttt{square\_log\_limit}}
in \texttt{RectangularCountEntropy}
gives a square limit for symmetric separately subadditive rectangular
log-counts with a linear area bound. Its scope is a square-limit
counterpart, not the weighted rectangular conclusions above.

\begin{thebibliography}{9}\small
\bibitem{a394697} T. Hildebrant, OEIS A394697, revision 24, 8 May 2026;
conjecture comment by V. Kot\v{e}\v{s}ovec, history revision 21, 6 May 2026.
\href{https://oeis.org/history?seq=A394697}{Entry history}.
\href{https://github.com/oeis/oeisdata/blob/77073cc8105b8042be07ec0d885da5ea5960ff3c/seq/A394/A394697.seq}{Pinned entry}.
\bibitem{a247413} V. Kot\v{e}\v{s}ovec, OEIS A247413, revision 41,
27 October 2026, formula identifying the A063443 limit.
\href{https://github.com/oeis/oeisdata/blob/77073cc8105b8042be07ec0d885da5ea5960ff3c/seq/A247/A247413.seq}{Pinned entry}.
\bibitem{a063443} R. Martin and contributors, OEIS A063443, revision
104, 5 November 2025; tiling definition and A. Woods's 2011 king-board
comment. \href{https://github.com/oeis/oeisdata/blob/77073cc8105b8042be07ec0d885da5ea5960ff3c/seq/A063/A063443.seq}{Pinned entry}.
\bibitem{kings} N. J. Calkin, K. James, S. Purvis, S. Race,
K. Schneider and M. Yancey, \emph{Counting Kings: As Easy As
$\lambda_1,\lambda_2,\lambda_3,\ldots$}, 18 July 2006, Sections 3--4.
\url{http://www.math.clemson.edu/~kevja/PAPERS/Kings2-2005REU.pdf}.
\bibitem{liang} K. Liang, \emph{Independent Set Enumeration in King
Graphs by Tensor Network Contractions}, arXiv:2505.12776v2,
v1: 19 May 2025, Proposition 3.1, pp. 9--10;
v2: 23 June 2025, Propositions 3.1 and 4.1, pp. 10--11 and 13--14.
\url{https://arxiv.org/abs/2505.12776v2}.
\bibitem{liang2} K. Liang, \emph{Independent Set Enumeration and
Estimation of Related Constants of Grid Graphs and Their Variants}, arXiv:2507.04007v2,
8 August 2025, Sections 3.2 and 4.
\url{https://arxiv.org/abs/2507.04007v2}.
\bibitem{hildebrant} T. Hildebrant, \emph{Transfer Operators and
Independence Polynomials for Strong Powers of Circulant Graphs},
arXiv:2604.14416v1, 15 April 2026, Proposition 2.5 and Remark 2.6.
\url{https://arxiv.org/abs/2604.14416v1}.
\bibitem{boundaries} R. H. Hardin, OEIS A067958, revision 25, and
V. Kot\v{e}\v{s}ovec, OEIS A212269, revision 34, especially the latter's
limit formula.
\href{https://github.com/oeis/oeisdata/blob/77073cc8105b8042be07ec0d885da5ea5960ff3c/seq/A067/A067958.seq}{A067958};
\href{https://github.com/oeis/oeisdata/blob/77073cc8105b8042be07ec0d885da5ea5960ff3c/seq/A212/A212269.seq}{A212269}.
\bibitem{capacity} V. Kot\v{e}\v{s}ovec and contributors, OEIS A189889,
revision 67, 11 April 2026.
\href{https://github.com/oeis/oeisdata/blob/77073cc8105b8042be07ec0d885da5ea5960ff3c/seq/A189/A189889.seq}{Pinned entry}.
\end{thebibliography}
\bigskip
\noindent\small\copyright\ 2026 Haobo Ma and Wenlin Zhang.
Licensed under Creative Commons Attribution 4.0 International
(CC BY 4.0): \url{https://creativecommons.org/licenses/by/4.0/}.
\end{document}
